% Nouveau EGA V (cote 25): archive cote 25, site page 84 (PDF page 85).
% typescript. Draft compared: none.
% AI-transcribed by Claude Opus 5.5 (claude-opus-5-5, adaptive thinking, effort high), Message Batches API, batch ega25b (msgbatch_01A7n5YuXubUPNK1zsepi8RG), 2026-11-13; not yet checked by eye.
\scanpage{25/84}
\mspage{\hand{78}}

14. \uline{Projections coniques}.

NB Nous avons déjà utilisé les projections coniques à diverses
\margin{$\|$}reprises, notamment à la fin du nº8, formulation de 10.4. et autres, et
le sorite qui suit devrait venir sans doute plutôt au début du paragraphe,
et éventuellement dans le paragraphe auxiliaire ``grassma\unsure{n}nien''.
Soit $C = P(F)$ une sous-variété linéaire de $P(E)=P$, de dimension
relative $r-m-1$ \xout{dans} sur $S$\inserted[typed]{i.e. de codimension $m+1$ dans $P$}, donc $F$ est un quotient de $E$, loc libre de
rang $r-m$ , $F=E/G$ , où $G$ est loc libre de rang $m+1$. On a défini
par voie algébrique au Chap II un morphisme
\[ p_C : P - C = P(E) - P(E/G) \longrightarrow P(G) , \]
que nous allons interpréter par voie géométrique, et qui s 'appelera
(en vertu d e la description qui suit) la\uline{projection conique} de centre $C$.
(NB On suppose $r-m-1$ compris entre $-1$ et $r-1$, i.e. $m$ compris entre $0$
et $r$, sans plus). Pour ceci, commençons par interpréter \xout{Q(G)xx} $P(G)$
comme un sous-\inserted{pré}schéma\inserted{fermé} de $\mathrm{Grass}^m(P)= \mathrm{Grass}_{r-m+1}(P)$, grâce à l'homomorphisme
de foncteurs évidents
\[ \xout{Gxx}\ P(G) \longrightarrow \mathrm{Grass}_{r-m+1}(E) \]
otebnu en considérant, pour chaque quotient inversible $G/G'$ de $G$, le
Module localement libre de rang $(r-m)+1$ $E/G'$ de $E$ (et kif kif après
tout changement de base). L'homomorphisme de foncteurs précédent est un
\uline{monomorphisme}, et comme le premier est propre sur $S$, le deuxième séparé,
c'est une immersion fermée. \xout{\unsure{(}}Plus généralement, il y aurait lieu d'expli-%
citer les immersions fermées des grassmaniennes de $G$, i/e. de $P(G)$,
dans celles de \inserted{$E$ i.e. de} $P=P(E)$. \xout{Ixix} L'image\inserted[typed]{(au sens foncteurs)} du morphisme obtenu est formé des
sous-variétés linéaires \xout{dexM} $L$\inserted[typed]{de la dimension voulue} de $P$ qui \uline{contiennent} $C$ . Désignant
par $Q(C)$ cette image dans le cas qui nous occupe (i.e. pour les dimensi-%
ons précisées plus haut), et identifiant $P(G)$ à $Q(C)$, le morphisme
de projection conique
\[ p_C : P - C \longrightarrow Q(C) \subset \mathrm{Grass}^m(P) \]

\mspage{\hand{84}}
